\section{The exact condition on a small relative category}

Fix a $\U$-small relative category $(M,\Sigma)$ and objects $a,b$. Let $L=L(M,\Sigma)$ be its infinity-categorical localization. Write
\[
 H_{a,b}(X)=\Hom_{M^X[\Sigma_X^{-1}]}(a_X,b_X),\qquad
 Z_{a,b}=\Map_L(a,b).
\]
The functor $M^X\to\Fun(NX,L)$ sends pointwise marked arrows to equivalences, and therefore induces a comparison
\[
 \alpha_X:M^X[\Sigma_X^{-1}]\longrightarrow\Ho\Fun(NX,L).
\]
Its map on constant-to-constant Hom-sets is written
$c_X:H_{a,b}(X)\to[BX,Z_{a,b}]$.

\begin{claim}{C.13}{Proved}{The comparison on constants is the specified path comparison}
For every cofinal path sequence $K_n=F_{t_n}(a,b)$, the path formula identifies $c_X$ with
\[
 \colim_n\pi_0\Fun(X,K_n)
 \longrightarrow[BX,Z_{a,b}],
\]
induced by evaluation and the maps $BK_n\to Z_{a,b}$ of C.8. This equality is natural in $X$, $a$, and $b$, and preserves composition.
\end{claim}
\begin{proof}
A representative on the left is a functor $X\to F_{t_n}(a,b)$, equivalently a functor $X\times I_{t_n}\to M$ with constant endpoints and marked backward edges. Applying $M\to L$ and inverting those edges gives the coherent family of morphisms from $a$ to $b$ represented by its image under $\alpha_X$. The construction of $Q$ in C.8 performs exactly these inversions on each interval, and applies to the same product diagram. Thus the two maps agree on every word. C.8 identifies the reparametrization and ladder relations and proves compatibility with concatenation, so they agree on the colimit and preserve composition. Naturality follows by restricting or whiskering the original product diagram.
\end{proof}

\begin{claim}{C.14}{Proved}{Equivalent exact criteria for invariance}
For the fixed pair $a,b$, the following conditions are equivalent.
\begin{enumerate}[label=(\roman*)]
\item $H_{a,b}$ inverts all nerve weak equivalences in $\Cat_\U$.
\item For every stage $n$, there are $m\geq n$ and a functor $r_n:\Pi_1K_n\to K_m$ such that
\[
 [r_n\eta_{K_n}]=[K_n\to K_m]\quad\text{in }\Cathot.
\]
\item $Z_{a,b}$ is a $1$-type, and for every $X$ the square of sets
\[
 \begin{array}{ccc}
 H_{a,b}(\Pi_1X)&\xrightarrow{c_{\Pi_1X}}&[B\Pi_1X,Z_{a,b}]\\
 \downarrow&&\downarrow\\
 H_{a,b}(X)&\xrightarrow{c_X}&[BX,Z_{a,b}]
 \end{array} \tag{C.14.1}
\]
is cartesian. The vertical maps are the specified restrictions along $\eta_X$.
\end{enumerate}
Applying these conditions to every pair of objects classifies the requested small relative categories. Condition (ii) is independent of the chosen cofinal sequence.
\end{claim}
\begin{proof}
The formula checked in C.20 and B.1 identify (i) with invertibility of the canonical ind-unit. In particular $H_{a,b}(\eta_{K_n})$ is surjective. Lift the element represented by $\id_{K_n}$ through that map. A representative of the lift is a functor $\Pi_1K_n\to K_j$, and equality of the resulting elements is witnessed at a common later stage. This is exactly (ii).

Conversely, assume (ii). For $f:X\to K_n$, the functor $r_n\Pi_1f:\Pi_1X\to K_m$ restricts to the given class of $f$ at a later stage. Hence $H(\eta_X)$ is surjective. To prove injectivity, represent two classes by $g,h:\Pi_1X\to K_n$, and move to a stage $j$ at which $g\eta_X$ and $h\eta_X$ are transformation-homotopic. Apply $\Pi_1$: the resulting natural isomorphism, and the inverse of $\Pi_1\eta_X$, give $\Pi_1g\cong\Pi_1h$. Postcompose with $r_j$ from (ii). Its specified homotopy identifies the two resulting functors with the transition images of $g,h$. Thus their classes agree at a further stage. This proves injectivity. Groupoids invert nerve equivalences, so (i) follows.

To obtain the mapping-space assertion from (i), use B.1's canonical identity $H(X)=\colim_n[BX,BK_n]$. For each $d\geq2$, choose a finite face-poset model $P_d$ of $S^d$. First-truncation invariance gives $H(P_d)=H(*)$. The homotopy-colimit equivalence of C.3 and R7 identifies $Z_{a,b}$ with the actual telescope of the $BK_n$. Compactness of $S^d$ gives
\[
 [S^d,Z_{a,b}]=\colim_n[S^d,BK_n]=\colim_n\pi_0BK_n=\pi_0Z_{a,b}.
\]
The equality is induced by the constant-map function. A nonzero based homotopy element cannot be carried to zero by the fundamental-group action; hence every higher homotopy group vanishes. This proof uses the actual telescope and compact spheres, rather than an unproved passage from ind-isomorphisms to colimits.

In (C.14.1), (i) makes the left vertical map bijective by B.1, and the proved $1$-type assertion makes the right vertical map bijective. A commuting square with both vertical maps bijective is cartesian. Conversely, under (iii) the right vertical map is bijective, and its pullback is the left map, which is therefore bijective. Factoring through $\Pi_1$ proves (i). This proves all directions and independence of the sequence.
\end{proof}

\begin{claim}{C.15}{Proved}{A finite relative example with a nontrivial mapping-space fundamental group}
For any small groupoid $G$, form $M$ by adjoining objects $a,b$, with no arrows between them, no arrows from $G$ to them, and one arrow from each of $a,b$ to every object of $G$. Mark all morphisms in the full cone $\{b\}*G$ and the identity of $a$. Then, for every small $X$,
\[
 H_{a,b}(X)\cong\pi_0\Fun(X,G),\qquad
 \Map_{L(M,\Sigma)}(a,b)\simeq BG. \tag{C.15.1}
\]
The comparison sends a functor $g:X\to G$ to the roof $a_X\to g\leftarrow b_X$. Every constant-endpoint presheaf $H_{x,y}$ of this relative category is invariant. In particular $G=\mathbf BC_2$ gives a finite relative category and a connected invariant mapping space with fundamental group $C_2$.
\end{claim}
\begin{proof}
An admissible path from $a$ to $b$ leaves $a$ once, on a positive edge into $G$. It first reaches $b$ on a negative edge from $b$ to $G$. Before that first arrival, its $G$-vertices form a nonempty interval. It never returns to $a$. For a path $c$, let $k(c)$ be its first $G$-vertex and $l(c)$ its first $b$-vertex. A vertical transformation $c\to d$ satisfies $k(d)\leq k(c)$ and $l(d)\geq l(c)$: a vertex of $G$ can map only into $G$, and $a$ has no arrow to $b$. Thus the first $G$-interval of $d$ contains that of $c$.

Define $E_\tau(c)=c(k(c))$. On a transformation, use its component at $k(c)$ followed by transport inside the first $G$-interval of $d$ back to $k(d)$. Transport means the ordered composite of path arrows and inverses according to traversal direction. Naturality on that interval proves the transport equation; this equation proves preservation of composition, so $E_\tau:F_\tau(a,b)\to G$ is a functor.

Let $U_\tau c$ equal $c$ before $l(c)$ and be constantly $b$ thereafter. On a transformation use its original components before that index and the unique arrows out of $b$ afterwards. Every square is either an original naturality square or follows from uniqueness. This gives a functor and a natural transformation $U_\tau\Rightarrow\id$.

Let $k_0$ be the vertex just after the first positive edge of $\tau$, and let $l_0$ be the vertex just after its last negative edge. An admissible type has $k_0<l_0$. Define $s_\tau:G\to F_\tau$ to have value $a$ before $k_0$, value $g$ on the middle interval, and value $b$ from $l_0$ on. The middle arrows are identities; the entrance and exit are the unique cone arrows. Then $E_\tau s_\tau=\id_G$. There is a natural transformation $U_\tau\Rightarrow s_\tau E_\tau$: use transport back to $c(k(c))$ on the original first $G$-interval and the unique arrows out of $a$ or $b$ on the additional middle vertices. The inequalities $k_0\leq k(c)<l(c)\leq l_0$ prove that these are all cases. The same transport equation verifies naturality.

We have constructed literal functors with
\[
 \id_{F_\tau}\ \longleftarrow\ U_\tau\
     \longrightarrow\ s_\tau E_\tau,
 \qquad E_\tau s_\tau=\id_G. \tag{C.15.2}
\]
Applying $\Fun(X,-)$ preserves these transformations for arbitrary $X$. A type pullback cannot skip the positive entrance and negative exit of the first $G$-interval in one edge, because a type edge traverses a run of one sign. The first intervals therefore intersect, and their specified groupoid transport gives $E_\tau u^*\cong E_\sigma$, compatibly with composition of type maps. Hence the transition maps of any cofinal sequence become the identity on $\pi_0\Fun(X,G)$. The path formula proves the first isomorphism of (C.15.1), including mixed-valued intermediate diagrams. The same coherent comparisons on classifying spaces identify the homotopy colimit with $BG$, and C.7--C.8 give the second equivalence with its specified localization map.

For completeness, consider all other endpoint pairs. Write $C=\{b\}*G$. Every $g\in G$ is isomorphic to $b$ in each localized diagram category, through the marked arrow $b_X\to g_X$. Thus it suffices to consider $a,b$. A pointwise admissible path from $a$ back to $a$ cannot leave $a$: a forward arrow cannot return there and a marked backward arrow incident to $a$ is its identity. Consequently $H_{a,a}(X)$ is a singleton. No pointwise path goes from $b$ to $a$, so $H_{b,a}(X)$ is empty for nonempty $X$ and a singleton for empty $X$. Paths from $b$ to $b$ stay inside $C^X$. All its arrows are marked and $b_X$ is initial, so its groupoid completion has a singleton endomorphism set at $b_X$. Hence $H_{b,b}(X)$ is a singleton. These descriptions, together with (C.15.1), prove invariance for every pair, since a nerve weak equivalence preserves emptiness. The finite category for $G=\mathbf BC_2$ has three objects and six arrows; the specified operations and path retractions are checked by the program \texttt{code/c2\_suspension\_checks.py}.

\end{proof}

\begin{claim}{C.16}{Proved}{The scope of full comparison functors}
Bijectivity of the constant-to-constant Hom comparisons $c_X$ for all $X$, together with $1$-truncated mapping spaces, is a sufficient instance of C.14. Equivalence of $\alpha_X$ for every $X$ entails these Hom-set bijections by restricting full faithfulness to the two constant objects. Cartesianness of (C.14.1) is the exact condition of C.14.
\end{claim}
\begin{proof}
Under the stated bijections, (C.14.1) has bijective horizontal maps and its right vertical map is bijective, so the square is cartesian. The proof of C.14 uses precisely those four sets and their maps. It imposes no condition on the remaining Hom-sets of $\alpha_X$. Thus the sufficient statement and the exact specified criterion follow without a diagram-rectification assumption.
\end{proof}
